%==========================================================
% BACKMATTER — Volume I: Zero to Bit
%==========================================================

\appendix

%----------------------------------------------------------
\chapter{Standard C and POSIX Headers}
\label{app:headers}
\addcontentsline{toc}{chapter}{Appendix A\quad Standard C and POSIX Headers}
\chaptermark{Standard Headers}
%----------------------------------------------------------

\begin{chapterintro}
The topics in this volume correspond to a well-defined set of standard C
and POSIX headers. This appendix maps those topics to the headers where
the real definitions live, so that moving from the text to actual code is
straightforward. Entries are grouped to match the order of the chapters.
\end{chapterintro}


\backmatter
\pagestyle{backmatter}

\chapter*{End of Volume\,I}
\addcontentsline{toc}{chapter}{End of Volume I}
\chaptermark{End of Volume I}

Volume~I has followed one thread: how information is encoded in digital
systems, starting at a voltage level and ending at compressed binary
streams. Integers, floating-point numbers, Unicode text, byte order,
aligned memory, pointers, error-detection codes, serialization formats ---
these are not separate topics. They are successive layers of agreed-upon
meaning stacked on top of the same physical signal.

Volume~II picks up where this one ends. The encodings you have read about
here exist in physical hardware --- transistors, logic gates, memory cells,
cache hierarchies --- that imposes its own constraints on how those encodings
are stored, moved, and transformed. Understanding those constraints is what
turns a working program into a fast one.

\medskip
\noindent\textit{If any chapter in this volume felt unclear, come back to it
after Volume~II. The hardware context makes many representational choices
obvious in retrospect.}

\clearpage